%%
%% This is file `main.tex' based on `sample-sigconf.tex' (q.v. for spurce of that,
%%
%% IMPORTANT NOTICE:
%% 
%% For the copyright see the original source file `sample-sigconf.tex'
%% in the `Sample' folder.
%%
%% For distribution of the original source see the terms
%% for copying and modification in the file samples.dtx.
%% 
%% This generated file may be distributed as long as the
%% original source files, as listed above, are part of the
%% same distribution. (The sources need not necessarily be
%% in the same archive or directory.)
%%
%% Commands for TeXCount
%TC:macro \cite [option:text,text]
%TC:macro \citep [option:text,text]
%TC:macro \citet [option:text,text]
%TC:envir table 0 1
%TC:envir table* 0 1
%TC:envir tabular [ignore] word
%TC:envir displaymath 0 word
%TC:envir math 0 word
%TC:envir comment 0 0
%%
%%
%% The first command in your LaTeX source must be the \documentclass command.

%% NOTE that a single column version is required for 
%% submission and peer review. This can be done by changing
%% the \doucmentclass[...]{acmart} in this template to 
%%\documentclass[manuscript,review,anonymous]{acmart}
%% This version is used for drafting and final submission
\documentclass[sigconf]{acmart}
\settopmatter{printacmref=false}


%% 
%% To ensure 100% compatibility, please check the white list of
%% approved LaTeX packages to be used with the Master Article Template at
%% https://www.acm.org/publications/taps/whitelist-of-latex-packages 
%% before creating your document. The white list page provides 
%% information on how to submit additional LaTeX packages for 
%% review and adoption.
%% Fonts used in the template cannot be substituted; margin 
%% adjustments are not allowed.

%%
%% \BibTeX command to typeset BibTeX logo in the docs
\AtBeginDocument{%
  \providecommand\BibTeX{{%
    \normalfont B\kern-0.5em{\scshape i\kern-0.25em b}\kern-0.8em\TeX}}}

%% The following are not a requirement, delete if not using
\acmDOI{}
\acmISBN{}
\setcopyright{none}
\acmConference[ACL]{}{Algorithms for Computational Logic}{04-10-2025}

\makeatletter
\patchcmd{\maketitle}{\@copyrightyear.}{}{}{}
\makeatother

%%
%% Submission ID.
%% Use this when submitting an article to a sponsored event. You'll
%% receive a unique submission ID from the organizers
%% of the event, and this ID should be used as the parameter to this command.
%%\acmSubmissionID{123-A56-BU3}

%%
%% For managing citations, it is recommended to use bibliography
%% files in BibTeX format.
%%
%% You can then either use BibTeX with the ACM-Reference-Format style,
%% or BibLaTeX with the acmnumeric or acmauthoryear sytles, that include
%% support for advanced citation of software artefact from the
%% biblatex-software package, also separately available on CTAN.
%%
%% Look at the sample-*-biblatex.tex files for templates showcasing
%% the biblatex styles.
%%

%%
%% The majority of ACM publications use numbered citations and
%% references.  The command \citestyle{authoryear} switches to the
%% "author year" style.
%%
%% If you are preparing content for an event
%% sponsored by ACM SIGGRAPH, you must use the "author year" style of
%% citations and references.
%% Uncommenting
%% the next command will enable that style.
%%\citestyle{acmauthoryear}

%%
%% end of the preamble, start of the body of the document source.

%% % Location of your graphics files for figures, here a sub-folder to the main project folder
\graphicspath{{./images/}} 


\begin{document}

%%
%% The "title" command has an optional parameter,
%% allowing the author to define a "short title" to be used in page headers.
\title{The Nurse Scheduling Problem}

%%
%% The "author" command and its associated commands are used to define
%% the authors and their affiliations.
%% Of note is the shared affiliation of the first two authors, and the
%% "authornote" and "authornotemark" commands
%% used to denote shared contribution to the research.
\author{Francisco Fonseca, 102492}
\authornote{Both authors contributed equally to this work.}
\email{franciscodiasfonseca@tecnico.ulisboa.pt}
\author{Gonçalo Rua, 102604}
\authornotemark[1]
\email{goncalonrua@tecnico.ulisboa.pt}
\affiliation{%
  \institution{Instituto Superior Técnico}
  \streetaddress{Av. Rovisco Pais}
  \city{Lisbon}
  \country{Portugal}
}

%%
%% By default, the full list of authors will be used in the page
%% headers. Often, this list is too long, and will overlap
%% other information printed in the page headers. This command allows
%% the author to define a more concise list
%% of authors' names for this purpose.
\renewcommand{\shortauthors}{Francisco Fonseca, Gonçalo Rua}

%%
%% This command processes the author and affiliation and title
%% information and builds the first part of the formatted document.
\maketitle

\section{Description of the Problem}
The problem in question is to create an optimal shift schedule for nurses working in a hospital during a given period of days, respecting some constraints:

\begin{enumerate}
    \item Each shift has a minimum number of nurses necessary for the normal functioning of the hospital.
    \item Each shift requires at least one nurse with managerial qualifications.
    \item A nurse's daily schedule has to be continuous, meaning that there are no gaps in between assigned shifts.
    \item Each nurse has a minimum and maximum daily workload.
    \item The total number of shifts that can be assigned to a nurse for the whole period is limited.
    \item A nurse cannot be assigned to shifts when it's known they're unavailable to work.
    \item All nurses' shift preferences, known beforehand, have to be taken into account, being an optimal solution one in which the most nurses are happy with their schedule.
\end{enumerate}

\section{Description of the Encoding}
This problem can be modeled as a MaxSAT problem. To solve it we first need to define variables which can then be used to encode the constraints above. This encoding can then be fed into a MaxSAT solver and an optimal solution will be generated. Let us define the variables of the problem.
\subsection{Definition of Variables}
\begin{itemize}
    \item Let $D$ be a positive integer representing the number of days.
    \item Let $S$ be a positive integer representing the number of shifts per day, all with the same duration, that cover the full 24 hours. For example if $S=4$, each shift will have 6 hours.
    \item Let $R_{s,d}$ define the minimum number of nurses required for shift $s$ $(1\leq s \leq S)$ at day $d$ $(1\leq d \leq D)$.
    \item Let $N = {\{N_1, N_2,...,N_n}\}$ be a set of $n$ positive integers representing the available nurses to be assigned to work in the shifts.
    \item Let $M \subseteq N$ define the subset of nurses with managerial qualifications.
    \item Let $k_{{min}}$ and $k_{{max}}$ define the minimum and maximum number of shifts for nurse $N_k$ $(1 \leq k \leq n)$ in any given day.
    \item Let $H_k$ define the maximum total number of shifts that can be assigned to nurse $N_k$ for the whole period.
    \item Let $P_{k,s,d}$ be a set of $n \times S \times D$ integers representing the cost of assigning nurse $N_k$ to shift $s$ in day $d$. If $P_{k,s,d} = -1$, nurse $N_k$ is not available for shift $s$ in day $d$. All other values in $P$ are positive.
    \item Let $X = {\{x_1,x_2,...,x_i \}}$, for any $i \leq n \times S \times D$, be a set of variables, each taking a value in $\{0,1\}$ such that $x_i=1$ iff the nurse is assigned to that shift. What nurse is assigned to what shift in which day is clarified below.
    \item Let there be within $X$ $k$ subsets $X_k$ of $S \times D$ variables, each associated with nurse $N_k$. For example, $X_1$ is associated with nurse $N_1$ and contains the first $S \times D$ variables, $X_2$ is associated with nurse $N_2$ and contains the next $S \times D$ variables, and so on. 
    \item Let there be within each $X_k$ $D$ subsets $X_{k,d}$ of $S$ variables, each associated with nurse $N_k$ and day $d$. For example, $X_{1,1}$ represents $d=1$ of $N_1$ and contains the first $S$ variables, $X_{1,2}$ represents $d=2$ of $N_1$ and contains the next $S$ variables, and so one.
    \item Within each $X_{k,d}$, there are $S$ variables, each representing the participation of nurse $N_k$ for day $d$ in each shift. For example, $X_{1,1,1}=x_1$ and encodes the participation of nurse $N_1$ in shift $s=1$ in day $d=1$, $X_{1,1,2}=x_2$ and encodes the participation of nurse $N_2$ in shift $s=1$ in day $d=1$, and so on. We can essentially represent $x_i$ as $X_{k,d,s}$ and vice-versa. An example will be given below for better understanding.
    \item Let $L_{d,s}$ be a subset of $X$ that contains all variables $x_i$ associated with shift $s$ of day $d$.
\end{itemize}

Let $D=2$, $S=2$ and $N=\{N_1,N_2,N_3 \}$ such that $X = \{x_1,x_2,...x_{12} \}$.

\vspace{0.3em}
\begin{center}
    \begin{tabular}{ccc}
        \toprule
        \textbf{Nurse 1} & \textbf{Shift 1} & \textbf{Shift 2} \\
        \midrule
        Day 1 & $x_1$ & $x_2$ \\
        Day 2 & $x_3$ & $x_4$ \\
        \bottomrule
    \end{tabular}

    \vspace{1em} % Adds a controlled vertical space between tables

    \begin{tabular}{ccc}
        \toprule
        \textbf{Nurse 2} & \textbf{Shift 1} & \textbf{Shift 2} \\
        \midrule
        Day 1 & $x_5$ & $x_6$ \\
        Day 2 & $x_7$ & $x_8$ \\
        \bottomrule
    \end{tabular}

    \vspace{1em} % Adds a controlled vertical space between tables

    \begin{tabular}{ccc}
        \toprule
        \textbf{Nurse 3} & \textbf{Shift 1} & \textbf{Shift 2} \\
        \midrule
        Day 1 & $x_9$ & $x_{10}$ \\
        Day 2 & $x_{11}$ & $x_{12}$ \\
        \bottomrule
    \end{tabular}
\end{center}
\vspace{0.3em}

Any $X_{k,s,d}$ can be represented as $x_i$:
\[
    i = (S \cdot D)(k - 1) + S(d - 1) + s
\]

%% ACHO QUE ISTO ESTÁ UM BOCADO CONFUSO E DESNECESSÁRIO, SQUE MUDA-SE!!!!!!!!!!!!!!!!!
On the contrary, for $i=7$, we know that $x_i \in X_2$ due to the size of any $X_k$ being $D*S$, and further we know that $x_i \in X_{2,2}$ due to the size of any $X_{k,d}$ being $S$. Finally, we can use the previous equation to find that $x_7 = X_{2,2,1}$.

Thus, $x_i = X_{k,s,d} = 1$ iff nurse $N_k$ is assigned to shift $s$ in day $d$. For instance, $x_9=1$ iff nurse $N_3$ is assigned to the shift $s=1,d=1$. Within our program this set X is cleanly managed by the PySAT IDPool class.

\subsection{Modeling Constraints}
Having defined the necessary variables, we can now describe the encoding we chose to model the initial constraints.

\begin{equation}
    \sum_{x_j \in L_{d,s}} x_j \geq R_{s,d}, \quad \forall d \in [1,D], s \in [1,S]
\end{equation}

\begin{equation}
    \sum_{x_j \in L_{s,d}} x_j \geq 1, \quad \forall k,s,d: s \in [1,S], d \in [1,D], N_k \in M
\end{equation}

\begin{equation}
    \begin{gathered}
        (\neg X_{k,s,d} \lor X_{k,s+1,d} \lor \neg X_{k,v,d}), \\
        \forall k,s,d,v: k \in [1,n], s \in [1,S-1], d \in [1,D], v \in [s+2,S]
    \end{gathered}
\end{equation}

\begin{equation}
    \begin{aligned}
        \sum_{x_j \in X_{k,d}} x_j &\geq k_{\min}, \quad \forall k \in [1,n], d \in [1,D] \\
        \sum_{x_j \in X_{k,d}} x_j &\leq k_{\max}, \quad \forall k \in [1,n], d \in [1,D]
    \end{aligned}
\end{equation}

\begin{equation}
    \sum_{x_j \in X_k} x_j \leq H_k, \quad \forall k \in [1,n]
\end{equation}

\begin{equation}
    (\neg X_{k,d,s}), \quad \forall k,d,s: P_{k,s,d} < 0
\end{equation}

\vspace{0.5em}

We guarantee an optimal solution will be found by adding the following soft clauses:
\begin{equation}
    (\neg X_{k,d,s}) \text{ with weight } P_{k,s,d}, \quad \forall k,s,d: P_{k,s,d} \geq 0
\end{equation}

\section{Algorithm and Configuration}
With the encoding defined and coded into Python, we fed the formula into the RC2 MaxSAT solver \cite{rc2} available in PySAT, using the CaDiCaL 1.9.5 oracle \cite{cadical} underneath. We chose $cadical195$ simply because it far outperformed every other option we tried. All cardinality constraints use a Modulo Totalizer encoding \cite{mtotalizer}.

\section{Installing and running}
To solve this project we used Python. The only dependency used in the project was PySAT which can be installed by running the following command: \verb|pip3 install pysat|. To run the program, simply execute:
\begin{verbatim}
python3 project1.py < input.txt > output.txt
\end{verbatim}
Note that \texttt{input.txt} and \texttt{output.txt} should be replaced with your actual file names.

%%
%% The next two lines define the bibliography style to be used, and
%% the bibliography file.
\bibliographystyle{ACM-Reference-Format}
\bibliography{references.bib}

%%
\end{document}
\endinput
%%
%% End of file `main.tex'.
